\subsection{Finite, noisy calibration of interaction-based tree selection}
\label{note:finite_calibration_morphology}

The exact-target construction leaves open whether a useful tree can be chosen from a limited sample. We therefore tested selection from finite noisy observations without access to the target's exact Fourier coefficients. The forward model was the eight-input algebraic tree used in the preceding representability analysis. Each of its seven internal nodes used the four-coefficient local computation in Eq.~\ref{eq:s_multiaffine_node}, giving 28 fitted coefficients, 14 edges and a root-only output. Leaves received individual named inputs; no global input mixer, additional decoder or bypass was available. This experiment concerns that constrained multilinear model, rather than conductance dynamics or biologically local plasticity.

Tasks were drawn from four equally weighted families: sums of four disjoint pair products, sums of two complementary quartic products, sums of nested products of degrees 2, 4, 6, 8, and a random-interaction control. The control selected four distinct nonconstant monomials uniformly from all degree-2–8 masks, independently of the candidate trees. Each task received a random input permutation where applicable, independent random coefficient signs and magnitudes sampled uniformly between 0.5 and 1.5. Coefficients were normalized to unit squared norm. Under independent Rademacher inputs, each target therefore had mean zero and population variance one. These random coefficients generated new functions; they did not preserve an identical input-gradient second moment across tasks. The earlier exact capacity experiment provides the separate isospectral comparison.

The primary candidate pool retained the same twelve trees: three shapes (balanced, comb and 3+5) under four fixed input assignments. An $\ell_1$-penalized linear estimator (Lasso) received input/output observations and the dictionary of all 255 nonconstant Walsh monomials. It received no family label, planted interaction support, true coefficient, teacher derivative or full-domain truth table. Its prior information was the known eight-variable Rademacher domain and a sparse multilinear dictionary. The intercept was fitted separately and excluded from the interaction score, because a global constant does not constrain subtree compatibility.

Calibration comprised 64, 256 or 1,024 independent input draws, with independent Gaussian label noise of standard deviation zero or 0.5. The primary condition was 256 queries with noise standard deviation 0.5. The first 75\% of observations were available for fitting the Lasso and short-pilot trees; the remaining 25\% formed a held-out calibration subset used to choose among the pilot fits (the calibration gate). Thus the primary estimator used 192 observations, while the shared total calibration-query budget was 256. Let $n_{\rm fit}$ be the number of fitting observations and $\operatorname{sd}(y_{\rm fit})$ their label standard deviation. The Lasso penalty was $c\,\operatorname{sd}(y_{\rm fit})\sqrt{2\log(256)/n_{\rm fit}}$. Fitting used cyclic coordinate descent, tolerance $10^{-8}$ and at most 20,000 iterations. The scale $c$ was selected from 0.15, 0.35, 0.65 using five development seed blocks. The frozen value was 0.15. The same development outcomes fixed the best constant candidate and a lookup rule assigning a candidate to the estimated input-gradient rank needed to explain 99.5\% of its energy. Both reduced to \texttt{balanced\_p0}; bins with fewer than two development tasks used the constant fallback.

For each candidate subtree, we matricized the estimated coefficients across the subtree/complement split, removed the row corresponding to the constant function inside the subtree, and computed the squared singular-value tail beyond its first component. The selector minimized the maximum normalized tail. Scores within $10^{-10}$ were treated as ties; ties were resolved by the sum of overlapping cut tails and then by lexicographic candidate name, using the same tolerance. The maximum is a population-MSE lower bound when computed from the true coefficients. Under noisy estimation it is a selection score, not a certified lower bound on the true task. The summed tails are a tie-breaking heuristic and are not a bound.

Baselines comprised fixed representatives of each shape, the development-best fixed candidate, the development-trained estimated-rank-only rule, uniform random choice averaged exactly over the candidate pool, and a two-sweep fitting pilot. The pilot independently fitted every candidate on the same calibration-fit observations from one fixed initialization and selected the smallest MSE on the calibration gate. Pilot parameters were discarded before final training, so the selected candidate received no additional training from this procedure. A secondary selector applied the previously audited subset dynamic program to the estimated coefficients at the primary calibration condition only. It could select a tree outside the twelve-candidate pool while retaining the same coefficient and edge budget.

All model definitions, estimator choices, baselines, training settings, primary comparisons and effect-size margin were frozen after development. Twenty new seed blocks generated one new function from each family, giving 80 confirmatory tasks. Each task was evaluated under all six calibration conditions. Cryptographic hashes recorded every calibration score, pilot choice, adaptive-tree choice and calibration dataset before any confirmatory candidate completed final training or exposed its final test endpoint. Two implementation errors stopped workers before fitting: lexicographic JSON ordering of numerical node IDs, and a launcher importing a historical module with the same name. The failure logs were retained. Separately hashed amendments corrected only node decoding and module dispatch. The original frozen selector files and all sealed choices remained unchanged.

For each task, final candidate training used a separate fixed stream of 1,024 input draws with Gaussian label noise of standard deviation 0.3. Every candidate used 32 alternating least-squares sweeps and two paired random initializations. Coordinate updates solved the four-parameter local least-squares subproblem, with output-preserving gauge normalization. A separate 512-observation noisy validation stream selected between the two final restarts. No test observation was used for fitting, restart choice or calibration. The short pilot and final fits used separate data and began independently from their declared initializations.

Final evaluation used 4,096 independent test inputs. The primary endpoint was MSE against the clean target, divided by its known population variance of one. MSE against independently noisy test labels and exact MSE over all 256 input patterns were secondary evaluation-only endpoints. Because the input domain is finite, independent streams can contain the same input pattern. They represent independent observations with fresh noise and distinct sample identifiers; they are not partitions of a reused noisy-label cache. The task generator knows the target coefficients to draw observations, but selectors receive only their stated calibration inputs. Complete coefficients and full-domain values are reserved for post-seal target-informed references and evaluation.

The twelve fixed candidates, the estimated adaptive tree and a privileged true-coefficient adaptive reference were each fitted twice per task, giving 2,240 confirmatory fits. The same final candidate outcomes were reused across calibration conditions. Primary regret was the selected candidate's independent-test normalized MSE minus the retrospective minimum among the twelve fixed candidates, each represented by its validation-selected restart. The adaptive selector's regret against this fixed pool can be negative; such values were retained and describe an advantage from searching a larger tree space. The true-coefficient cut selector, true-coefficient adaptive tree and retrospective best trained candidate in the pool were explicitly privileged references, not deployable baselines.

Uncertainty used 20 seed blocks. The four family-specific paired differences were averaged within each seed before bootstrapping 10,000 seed samples. The two primary comparisons were the estimated-cut selector against the development-best fixed candidate and against the two-sweep pilot, at 256 queries/noise 0.5. Descriptive intervals used 95\% coverage. Bonferroni-adjusted two-sided 97.5\% intervals were used for the two primary comparisons. Meaningful superiority required both a positive adjusted lower bound and mean improvement of at least 0.01 normalized MSE. Other calibration sizes, the noise-free conditions, family-specific results and adaptive construction were prespecified secondary analyses. Failed fits and convergence warnings were retained; nonfinite or linear-algebra failures received the predeclared loss of 1e6 and were counted explicitly.

The prespecified task mixture defines the scope of inference. Confirmation tests fresh functions and observations within that mixture; it does not establish transfer to unknown task families, a scalable compiler at larger input counts, or optimal biological dendritic morphology. The fixed Walsh dictionary grows exponentially with the number of inputs, and the tree fitter is an alternating least-squares procedure rather than a synapse-local plasticity rule.

Supplementary Fig.~\ref{fig:supp_morphology_calibration} reports the complete
calibration sensitivity, family-specific effects, primary adjusted contrasts
and implementation cost. Section~\ref{note:morphology_end_to_end} tests
subsequent gradient learning on estimated architectures. The separate study
in Section~\ref{note:morphology_credit_learning} uses oracle-compatible trees
to examine task, assignment and feedback-resolution effects.
